\documentclass[11pt]{article}
\usepackage{fontspec,amsmath,amssymb,amsthm,hyperref}
\usepackage[margin=1in]{geometry}
\setlength{\emergencystretch}{3em}
\title{Constructed iid particle comparison with uniform timewise L\ensuremath{{}^{2}} control}
\author{LangevinLean mathematical exposition}
\date{}
\begin{document}\maketitle
\noindent Source module: \texttt{ProximalLangevin/MeanFieldParticles.lean}. Compile this proof document with XeLaTeX or LuaLaTeX. The accompanying HTML pairs the mathematical argument with the exact Lean statements and records the scope of the formalization.


For a potential with globally \(L\)-Lipschitz gradient, compare a constructed mean-field solution with actual finite-particle coordinate solutions on the iid product source. The coefficient is chosen before all particle counts and the norm is the supremum of timewise \(L^2\) errors.


\section{The actual centered force is continuous and bounded on a finite horizon}


Define \(\mathcal C_V(X)=P_\theta\nabla V(X)-\operatorname{const}\mathbb E[P_\theta\nabla V(X)]\) using bounded linear centering and the actual Lipschitz drift. This composition is continuous. For continuous \(Y_t\), the image of compact \([0,T]\) is compact; its centered-force norm therefore has a finite bound \(C\ge0\). This \(C\) depends on the reference and horizon and is chosen independently of \(N\).




\textbf{Correspondence with the displayed declarations.}


\begin{itemize}

\item \texttt{centeredMeanFieldForce}: The centered actual parameter force of one mean-field copy.
\item \texttt{exists\_centeredMeanFieldForce\_bound}: Continuity of the constructed reference supplies a finite force coefficient on a compact horizon. This coefficient is independent of particle count.
\end{itemize}


\noindent Lean declarations: \texttt{centeredMeanFieldForce}, \texttt{centeredMeanFieldForce\_continuous}, \texttt{exists\_centeredMeanFieldForce\_bound}.

\section{Copies inherit the actual equation by bounded linearity}


The copy map \(\mathcal R_N\) is bounded linear, so it commutes with the Bochner time integral and preserves differences. Apply it to the mean-field equation to get the exact copied-reference forcing. Define

\[C_T=2\big[(C+2Lb)\Psi(2L,T)+b\big].\]

Here \(C\) bounds the actual centered force and \(b/\sqrt N\) bounds the extra parameter input. The factor two converts the coordinate maximum norm into the sum of parameter and one latent error.




\textbf{Correspondence with the displayed declarations.}


\begin{itemize}

\item \texttt{iidMeanFieldCopies\_integral\_equation}: The actual iid reference inherits its integral equation by a bounded linear map, so its forcing is constructed from the mean-field solution.
\item \texttt{meanFieldParticleConstant}: Finite-horizon coefficient in the tagged-coordinate L\ensuremath{{}^{2}} comparison. \texttt{C} bounds the actual centered reference force, and \texttt{b} bounds the parameter Brownian input after multiplication by \texttt{sqrt N}.
\end{itemize}


\noindent Lean declarations: \texttt{iidMeanFieldCopies\_integral\_equation}, \texttt{meanFieldParticleConstant}.

\section{Construct the particle solution and use the exact iid residual}


For each \(N>0\), Picard constructs the actual particle solution \(Z\) with copied latent input and any continuous extra parameter input \(A\) bounded by \(b/\sqrt N\). The copied mean-field reference has drift residual norm exactly \(\|\mathcal C_V(Y_t)\|_2/\sqrt N\le C/\sqrt N\). Apply the forcing-comparison theorem to the \(2L\)-Lipschitz particle drift. Both tagged coordinate errors are at most the product-state norm, hence

\[\|\Theta_t^N-\theta_t\|_2+\|X_t^{i,N}-X_t^i\|_2\le C_T/\sqrt N,\qquad0\le t\le T.\]

The theorem retains the constructed integral equation and exact iid initial state. It does not assume the final rate.




Continuity of the single-copy centered force along the constructed mean-field path implies a finite bound \(C\) on the compact interval \([0,T]\). Construct the actual particle process by Picard with copied latent input and the extra parameter input \(A_N\). The copied reference equation follows by commuting its bounded linear copy map through the time integral. Its drift defect is at most \(C/\sqrt N\), by the exact iid residual theorem, and \(\|A_N\|\le b/\sqrt N\).



The scaled integral-forcing theorem, with particle Lipschitz constant \(2L\), yields \[\|Z_t-R_t\|\le\frac{\operatorname{gronwallBound}(0,2L,C+2Lb,T)+b}{\sqrt N}.\] Each parameter or tagged-latent difference is bounded by the maximum product norm; their sum is therefore at most twice this expression. This accounts for the factor two in \texttt{meanFieldParticleConstant}. The reference and \(C\) are chosen before \(N\), making the coefficient independent of particle count.




\textbf{Correspondence with the displayed declarations.}


\begin{itemize}

\item \texttt{exists\_particle\_comparison\_of\_meanField\_solution}: The constructed particle solution is uniformly close, in each coordinate L\ensuremath{{}^{2}} norm, to the genuine iid copies of an actual mean-field integral solution. The reference force coefficient is finite and independent of \texttt{N}.
\end{itemize}


\noindent Lean declarations: \texttt{exists\_particle\_comparison\_of\_meanField\_solution}.

\section{Construct both systems with actual Brownian inputs}


The actual centered reference force is a continuous L\ensuremath{{}^{2}}-valued function of time. Compactness of \([0,T]\) supplies a finite \(C\ge0\) such that


\[
\sup_{0\le t\le T}\|
 \nabla_\theta V(Y_t)-\mathbb E_Q\nabla_\theta V(Y_t)\|_2\le C.
\]


This \(C\) is chosen \textbf{before the quantifier over all particle counts}; it depends on the reference potential, initial law and horizon, not on \(N\). Subtract the common latent input and the extra parameter Brownian input in the actual integral equations. The proof applies Grönwall to the remaining drift integral, so it never differentiates Brownian noise. Define


\[
\Psi(a,T)=\begin{cases}(e^{aT}-1)/a,&a>0,\\T,&a=0,\end{cases}
\qquad
C_T=2\big[(C+2Lb)\Psi(2L,T)+b\big].
\]


For every \(N>0\), the theorem constructs the actual particle integral solution with the prescribed iid initialization and proves, for every tagged particle \(i\),


\[
\sup_{0\le t\le T}\left(
 \|\Theta_t^N-\theta_t\|_2+\|X_t^{i,N}-X_t^i\|_2\right)
 \le\frac{C_T}{\sqrt N}.
\]


For the actual split-smoothed potential \(V=g_1+g_2^\eta\), use the derived constant \(L=L_1+1/\eta\). The coefficient can depend on \(\eta\); uniformity in smoothing is not claimed. This is a constructed iid propagation-of-chaos replacement in \textbf{uniform timewise L\ensuremath{{}^{2}} norm}. It does not move the time supremum inside the expectation, as in printed Proposition C.2; the expected path-supremum norm is stronger; the proved uniform timewise L\ensuremath{{}^{2}} estimate does not establish it.




\textbf{Correspondence with the displayed declarations.}


\begin{itemize}

\item \texttt{exists\_brownian\_meanField\_particle\_comparison}: Both the mean-field reference and finite-particle coordinate processes are constructed from actual Brownian-coordinate inputs. The parameter-noise \texttt{N-¹/\ensuremath{{}^{2}}} bound is proved from those laws by \texttt{empiricalParameterNoise\_bound}. The conclusion controls each time's L\ensuremath{{}^{2}} norm uniformly over the horizon.
\end{itemize}


\noindent Lean declarations: \texttt{exists\_brownian\_meanField\_particle\_comparison}.

\end{document}
